\documentclass[11pt]{article}
\usepackage[margin=1in]{geometry}
\usepackage{amsmath,amssymb,amsthm}
\usepackage{booktabs}
\usepackage{hyperref}
\usepackage{times}

\title{Learned Latent Curves: A Closed-Loop Framework for\\ Corpus Embedding and Recursive Skill Auditing}
\author{Shant Tchatalbachian \\ \small yubi-OS}
\date{August 4, 2026}

\begin{document}
\maketitle

\begin{abstract}
We describe a family of three composed techniques used to audit and prioritize
improvement work across a corpus of engineering artifacts (software skills,
self-documentation files): (1) a \emph{learned latent curve}, which re-expands
a single 1-D ordering coordinate $t$ into a fixed-width $D$-dimensional
embedding via a Fourier basis with learned frequencies; (2) \emph{curve-guided
recursive self-improvement} (curve-guided RSI), which uses sparse cells of the
fitted curve as a prioritization lens for gap-mapping and bounded
fixpoint-style editing cycles; and (3) a self-documentation variant of the
same pipeline, retargeted at heterogeneous memory-file corpora via a
per-corpus primitive basis and an explicit item-granularity rule. We give the
governing equations for each stage, the empirical validation obtained on a
corpus of software skills ($N=62$–$211$) and a self-documentation corpus
($N=17$–$154$), and we note the technique's distinction from an unrelated
prior work on weight-space reparameterization for neural network optimization
(arXiv:2607.09967) that shares only surface vocabulary (``curve'',
``curvature'') and no derivation.
\end{abstract}

\section{Introduction}

Many corpus-auditing problems reduce to the same shape: $N$ items live in a
high-dimensional feature space, and a practitioner wants (a) a single scalar
ordering of the items and (b) a fixed-width representation indexed by that
ordering, so that interpolation and sparse-region detection are well defined.
Dimensionality reduction (collapsing features to $t$) and dimensionality
\emph{expansion} (re-expanding $t$ into an embedding) are usually treated as
separate problems. This paper documents a single closed-form round trip
between them and its use as the fitting stage of a larger closed-loop audit
pipeline.

\section{The Learned Latent Curve}
\label{sec:curve}

For output dimension $j = 1, \dots, D$ (canonically $D = 384$) and 1-D
coordinate $t \in [0, 1]$, the model is

\begin{equation}
z_j(t) \;=\; a_{j,0} \;+\; \sum_{m=1}^{k} \Big( a_{j,m}\, \sin(2\pi f_m t)
\;+\; b_{j,m}\, \cos(2\pi f_m t) \Big),
\label{eq:model}
\end{equation}

with $k$ shared learned frequencies $f_1, \dots, f_k$ and per-output
coefficients $a_{j,m}, b_{j,m}$. Stacked over $j$, this is a curve
$\gamma : \mathbb{R} \to \mathbb{R}^D$. Writing the design vector

\begin{equation}
\phi(t) = \big[\, 1,\; \sin(2\pi f_1 t),\, \cos(2\pi f_1 t),\, \dots,\,
\sin(2\pi f_k t),\, \cos(2\pi f_k t) \,\big] \in \mathbb{R}^{1+2k},
\label{eq:design}
\end{equation}

the model is $\gamma(t) = C\, \phi(t)$ with coefficient matrix
$C \in \mathbb{R}^{D \times (1+2k)}$, giving a parameter count of

\begin{equation}
P \;=\; k \;+\; D(1 + 2k).
\label{eq:paramcount}
\end{equation}

At $D = 384$, $k = 8$: $P = 8 + 384 \times 17 = 6{,}536$ parameters, compared
against $N \times D$ target scalars.

\subsection{Closed-form coefficient solve}

With the frequencies $f$ held fixed, $C$ has a closed-form ridge solution

\begin{equation}
C^\star \;=\; Z^\top \Phi \, \big( \Phi^\top \Phi + \lambda I \big)^{-1},
\label{eq:ridge}
\end{equation}

where $\Phi \in \mathbb{R}^{N \times (1+2k)}$ is the design matrix and
$Z \in \mathbb{R}^{N \times D}$ is the target matrix. Only the $k$
frequencies require non-convex search; the closed-form ridge fit at any
candidate frequency vector is a sanity floor for the gradient-descent fit.

\subsection{Frequency parameterization}

Frequencies are stored in an unconstrained form and mapped through a softplus
to guarantee strict positivity and finite gradients near zero:

\begin{equation}
f_m \;=\; \operatorname{softplus}(\tilde f_m) \;=\; \log\!\big(1 + e^{\tilde f_m}\big), \qquad f_m > 0.
\label{eq:softplus}
\end{equation}

\subsection{Losses}

Let $t \in \mathbb{R}^N$, $Z \in \mathbb{R}^{N \times D}$ the targets, and
$\hat Z = \gamma(t)$ the fitted curve evaluated at the training coordinates.
The total objective is

\begin{equation}
\mathcal{L} \;=\; \mathcal{L}_{\mathrm{rec}} \;+\; \mathcal{L}_{\mathrm{freq}}
\;+\; \mathcal{L}_{\mathrm{smooth}} \;\big(+\, \mathcal{L}_{\mathrm{orth}}\big),
\label{eq:totalloss}
\end{equation}

\begin{align}
\mathcal{L}_{\mathrm{rec}} &= \operatorname{mean}\big( (\hat Z - Z)^2 \big), \label{eq:lrec}\\
\mathcal{L}_{\mathrm{freq}} &= \lambda_f \cdot \operatorname{mean}(f^2), \label{eq:lfreq}\\
\mathcal{L}_{\mathrm{smooth}} &= \lambda_s \cdot \operatorname{mean}
\Big( \big(\gamma(t_{\mathrm{grid}})[2:] - 2\gamma(t_{\mathrm{grid}})[1{:}\!-\!1]
+ \gamma(t_{\mathrm{grid}})[:\!-\!2]\big)^2 \Big), \label{eq:lsmooth}
\end{align}

where $t_{\mathrm{grid}}$ is a dense grid of $\sim$512 points spanning the
coordinate range and $\lambda_f \approx 10^{-4}$, $\lambda_s \approx 10^{-3}$.
$\mathcal{L}_{\mathrm{smooth}}$ regularizes the curve \emph{between} observed
points, where reconstruction loss alone is blind.

\subsection{Obtaining the coordinate $t$}

The default pipeline sets $t$ from the top principal component of a
$z$-scored $N \times D_{\mathrm{feat}}$ feature matrix, mapped to $[0,1]$; a
rank-uniformized variant replaces raw PC1 scores by their ranks to improve
design-matrix conditioning. The explained-variance ratio of PC1 (or
PC1+PC2 for a 2-D surface variant) is recorded as an explicit go/no-go gate,
with $\mathrm{PC1} \geq 0.40$ treated as evidence of usable low-rank
structure.

\section{Curve-Guided Recursive Self-Improvement}
\label{sec:rsi}

Curve-guided RSI composes the fitted curve with a gap-mapping and editing
loop. Coordinates are reduced to a 2-D surface $(u, v) \in [0,1]^2$ via
PC1+PC2 of a binary primitive-coverage matrix $C \in \{0,1\}^{N \times 9}$
lifted to $D = 384$ by a fixed seeded orthonormal projection $Q$:

\begin{equation}
Z \;=\; C \, Q^\top, \qquad (u_i, v_i) \;=\; \mathrm{PCA}_{1,2}(Z)_i.
\label{eq:liftpca}
\end{equation}

The $[0,1]^2$ plane is discretized into a $0.05 \times 0.05$ grid
($21 \times 21 = 441$ cells). For a cell $c = (u, v)$,

\begin{align}
\mathrm{neighbors}(c) &:= \big\{\, i \in \mathrm{corpus} : \| (u_i, v_i) - c \|_\infty \le r \,\big\}, \label{eq:neighbors}\\
\mathrm{is\_sparse}(c) &:= \big| \mathrm{neighbors}(c) \big| = 0, \label{eq:sparse}
\end{align}

with default radius $r = 0.05$. Cells with zero neighbors are gap
candidates; for each, a fresh-context gap-mapping subagent (focused on that
single item, never the whole corpus) is dispatched, and any flagged
``Extend'' gap is closed via a bounded editing loop capped at three cycles
per item. After all cycles, Stage 1 is re-run and the closed-loop claim is
verified by comparing the sparse-cell count before and after:

\begin{equation}
\Delta_{\mathrm{sparse}} \;=\; |\mathrm{sparse\_cells}|_{\mathrm{post}} - |\mathrm{sparse\_cells}|_{\mathrm{pre}} \;<\; 0
\quad\Longrightarrow\quad \text{gaps closed.}
\label{eq:delta}
\end{equation}

\section{Self-Documentation Variant}
\label{sec:self}

The self-documentation offshoot retargets Equations~\ref{eq:liftpca}–\ref{eq:delta}
at heterogeneous memory-file corpora (\texttt{SELF.md},
\texttt{SELF-CHANGELOG.md}, and eight additional preference/rule files),
under the granularity rule that \emph{each version is one corpus item}
(a changelog entry, a memory-file section, or a named row). Each corpus
receives its own 9-D binary primitive basis rather than sharing one basis
across structurally distinct document types, since a single shared basis was
found to flatten the per-corpus signal. Near-constant columns (coverage
$> 0.90$ or $< 0.10$) are dropped before the PCA lift in
Equation~\ref{eq:liftpca}.

\section{Empirical Validation}
\label{sec:results}

\begin{table}[h]
\centering
\begin{tabular}{@{}lrrrl@{}}
\toprule
Corpus / fit & $N$ & PC1(+PC2) & Holdout $R^2$ & Status \\
\midrule
Skill corpus, v1 (raw-content target) & 62 & 0.40 (1-D) & $-0.155$ & fail \\
Skill corpus, v2 (primitive-coverage target) & 213 & 0.40 (1-D) & $+0.183$ & pass \\
Skill corpus, v3 (2-D surface) & 211 & 0.4655 & $+0.4655$ & headline pass \\
Skill corpus, v4 (sentence-transformer target) & 211 & 0.0955 & $-0.005$ to $+0.130$ & fail (wrong target) \\
Curve-guided RSI, cycle 5 (69 skills) & 69 & 0.4615 & $+0.2244$ & pass \\
SELF-CHANGELOG.md, canonical & 17 & 0.9164 & $+0.9917$ & pass \\
SELF.md, canonical & 51 & 0.6952 & $+0.7041$ & pass \\
Combined memory-file corpus, v4 & 154 & 0.6479 & $+0.7877$ & pass \\
\bottomrule
\end{tabular}
\caption{Fit quality across corpora. PC1(+PC2) is the explained-variance
gate (Section~\ref{sec:curve}); Holdout $R^2$ is measured against held-out
items refit from their $t$ coordinate alone.}
\label{tab:results}
\end{table}

\section{Relation to Prior Work}

Equation~\ref{eq:model} shares no derivation with the weight-space
reparameterization proposed in arXiv:2607.09967 (``Learning in Curved Weight
Space: Exponential-Linear Weight Reparameterization for Improved
Optimization''), which maps a single raw neural-network weight scalar $w$
through a sign-aware symmetric-exponential/linear pathway pair to obtain an
effective weight $w_{\mathrm{eff}}$ for optimizer training, and is validated
by training-step reduction on autoregressive language models. The present
work instead maps a per-item ordering coordinate $t$ through a Fourier
basis to obtain a fixed-width corpus embedding, and is validated by holdout
reconstruction accuracy on document corpora. The overlap is limited to the
word ``curve''/``curvature'' and the presence of learnable shape parameters;
neither the domain, the basis function family, nor the fitting objective is
shared. The two are unrelated by construction, not merely by citation
omission --- a full side-by-side comparison is documented separately.

Within the embedding literature, Equation~\ref{eq:model} is a fitted
generalization of two established fixed-basis constructions: random Fourier
features \citep{rahimi2007} use \emph{fixed random} frequencies as a kernel
approximation, and transformer positional encodings \citep{tancik2020} use
\emph{fixed geometric} frequencies with no learned coefficients. The present
model instead learns both the frequencies and the coefficients from the
target embeddings directly.

\section{Conclusion}

The learned latent curve gives a closed-form, auditable, interpolatable
representation for small-$N$ corpora with a low-rank ordering structure.
Composed with sparse-cell detection and a bounded gap-closing loop, it
converts a one-shot corpus audit into a measurable, repeatable improvement
process with an explicit before/after metric (Equation~\ref{eq:delta}). The
same pipeline transfers, with corpus-specific primitive bases, from a
software-skill corpus to a self-documentation corpus.

\begin{thebibliography}{9}

\bibitem{rahimi2007}
A.~Rahimi and B.~Recht.
\emph{Random Features for Large-Scale Kernel Machines}.
Advances in Neural Information Processing Systems, 2007.

\bibitem{tancik2020}
M.~Tancik et al.
\emph{Fourier Features Let Networks Learn High Frequency Functions in Low
Dimensional Domains}.
Advances in Neural Information Processing Systems, 2020.

\bibitem{smith2026}
E.~Smith.
\emph{Learning in Curved Weight Space: Exponential-Linear Weight
Reparameterization for Improved Optimization}.
arXiv:2607.09967 [cs.LG], 2026.

\end{thebibliography}

\end{document}
